%-----------------------------------------------------------------------
\subsection{The Phase Problem (Crystallography)}
\label{sec:instance-phase-problem}
%-----------------------------------------------------------------------

The \emph{phase problem} is the central challenge of X-ray crystallography:
diffraction experiments measure the squared magnitudes $|F(\mathbf{k})|^2$ of
the structure factors, but the phases are lost.  Multiple electron density
distributions can produce identical diffraction patterns.
\citet{hauptman1953solution} developed direct methods that partially recover
the phases---work recognized with the 1985 Nobel Prize in Chemistry
\citep{karle1956structure}---but the fundamental ambiguity remains.

\paragraph{The explanation system.}
We formalize a minimal witness: a 2-element discrete signal.
\emph{Note:} For $N = 2$, the phase ambiguity is limited to permutation and
sign (the set $\{(a,b), (b,a), (-a,-b), \ldots\}$ sharing $a^2+b^2$).
For general $N$, the phase ambiguity grows combinatorially, which is why the
phase problem is hard in practice.  Our 2-element witness establishes that the
Rashomon property holds; the full crystallographic problem is strictly harder.
Define $\mathcal{S}_{\mathrm{phase}} = (\Params, \sH, Y, \mathsf{obs},
\mathsf{exp}, \incomp)$ as follows:
\begin{itemize}
  \item $\Params = \sH = \texttt{Signal2}$ (a pair of integers $(a, b)$
    representing the signal);
  \item $Y = \mathbb{Z}$ (the \emph{energy}: $a^2 + b^2$, modeling $|F|^2$);
  \item $\mathsf{obs} = \texttt{energy}$;
  \item $\mathsf{exp} = \mathrm{id}$ (the signal itself is the ``explanation''
    of the diffraction data);
  \item $\incomp(s_1, s_2) \iff s_1 \ne s_2$.
\end{itemize}

\paragraph{The Rashomon property.}
The witnesses are $s_1 = (1, 0)$ and $s_2 = (0, 1)$, which satisfy
$\texttt{energy}(s_1) = 1^2 + 0^2 = 1 = 0^2 + 1^2 = \texttt{energy}(s_2)$
and $s_1 \ne s_2$.  The Rashomon property is constructively witnessed with
zero axioms.

\paragraph{The impossibility.}
By Theorem~\ref{thm:universal}, no method of determining a signal from its
energy (diffraction intensity) can simultaneously be faithful (report the
actual signal), stable (assign the same signal to all inputs with the same
energy), and decisive (distinguish every pair of distinct signals).
In crystallographic terms: the magnitude of the structure factor does not
uniquely determine the electron density.

\paragraph{Empirical illustration.}
Gerchberg--Saxton phase retrieval from magnitude-only data produces 20
reconstructions with mean pairwise RMSD of 1.32--1.44 for unconstrained signals
(lengths 16--128), compared to 0.75--0.92 for positive signals where the
positivity constraint reduces phase ambiguity (ratio 1.5--1.8$\times$).
Statistical significance is assessed using the $N = 20$ per-reconstruction
mean RMSDs as independent units (not the 190 pairwise RMSDs, which are
pseudoreplicated): Mann--Whitney $U$, $p = 3.4 \times 10^{-8}$ per signal
length.

\paragraph{Resolution: Patterson maps.}
The classical partial resolution is the \emph{Patterson map}---the
autocorrelation function of the electron density, which is computable
directly from $|F(\mathbf{k})|^2$ without phase information.  The Patterson
map is the $G$-invariant observable: the symmetry group of phase rotations
acts on the signal space, and the autocorrelation is invariant under this
action.  Modern direct methods \citep{hauptman1953solution,karle1956structure} go further by exploiting
probabilistic constraints among phases, but even these methods cannot
eliminate the fundamental ambiguity without additional experimental data
(e.g., anomalous scattering or isomorphous replacement).

\paragraph{Lean formalization.}
The Lean~4 types are \texttt{Signal2} (a structure with fields
\texttt{a b : Int}), \texttt{energy} (the observation map
$s \mapsto s.a \cdot s.a + s.b \cdot s.b$), and
\texttt{phaseSystem : ExplanationSystem Signal2 Signal2 Int}.
The impossibility is \texttt{phase\_impossibility} in
\texttt{PhaseProblem.lean}, derived with zero axioms.
